% !TEX root = ../main.tex

\section{Tables and Imported Data}
\label{sec:tables-data}

The earlier tables show ordinary LaTeX source. NeoSetzer also provides an
\textbf{Insert Table} dialog for creating a grid, selecting alignment and
borders, merging constrained cells, and editing cell contents before inserting
LaTeX into the document.

\subsection{CSV and TSV practice data}
\label{sec:tables-data:csv}

This project includes \texttt{data/example-table.csv}. Open it in a text
editor, copy its rows, then open NeoSetzer's \menu{Insert Table} dialog and
choose \menu{Paste TSV/CSV}. The dialog detects tab-separated text first and
otherwise treats comma-separated CSV as the default; it also provides explicit
CSV delimiter choices and an \menu{Import CSV/TSV File} action.

The CSV file is deliberately a \emph{practice input} for the table generator.
It is not an instruction for LaTeX to read CSV automatically. After the dialog
inserts the resulting tabular code, build this project to see the table in the
PDF and adjust the generated source as you would in your own document.

\subsection{Long tables and cell structure}
\label{sec:tables-data:longtable}

For a table that may span pages, choose the \texttt{longtable} environment in
the generator or adapt the generated source. The existing table examples use
\texttt{booktabs}; they demonstrate professional horizontal rules and a
\texttt{\textbackslash multicolumn} total row. The editor's table workflow is
meant to speed up correct starting code, not to hide the LaTeX source from you.

\subsection{A suggested experiment}
\label{sec:tables-data:experiment}

Import \texttt{data/example-table.csv}, insert the result below this paragraph,
and add a caption and label. Then refer to it from a later paragraph with
\texttt{\textbackslash ref}. This links the table generator, labels,
completion, PDF preview, and Document Structure into one small reproducible
exercise.
